% !TeX root = main.tex
\section{Conclusiones}
\label{sec:conclusiones}

La implementación de segmentación interactiva basada en GrabCut permitió
obtener una región compatible con la referencia disponible para el corte
\texttt{VS-SEG-001}, a partir de una inicialización rectangular. Tras cinco
iteraciones, se registraron 47 falsos positivos y 84 falsos negativos. El
error cuadrático medio fue de $0{,}000500$ sobre la imagen completa y de
$0{,}026422$ dentro del rectángulo. Frente a la predicción que asigna fondo
a todos los píxeles, las discrepancias se redujeron aproximadamente un
$92{,}25$~\% en este caso.

Las etiquetas permanecieron sin cambios desde la tercera iteración hasta
la quinta, aunque los modelos de intensidad continuaron actualizándose.
El procedimiento combina así el ajuste de los modelos de intensidad con
la minimización de la energía mediante cortes de grafos.

El predominio del fondo hace que el error global sea pequeño incluso
cuando se omite toda la lesión. Por ello, su interpretación requiere
considerar también el error dentro del rectángulo, los falsos positivos y
negativos y su distribución espacial. En el caso evaluado, las
discrepancias se concentraron en el límite de la región segmentada.

El caso estudiado ilustra la aplicación de GrabCut a la segmentación
interactiva de una lesión en resonancia magnética. La inicialización
rectangular delimita la región de búsqueda, los modelos gaussianos
describen las intensidades de ambas clases y los enlaces entre vecinos
favorecen la coherencia espacial de la región obtenida.
